\subsection{Purpose}

This Software Requirements Specification (SRS) defines the requirements of
Dikaios Amanah, a secure, transparent, and ethically governed web-based
platform designed to connect Small and Medium Enterprises (SMEs) with
potential investors through a structured investment process.

The purpose of this document is to establish a clear and detailed
specification of the system's expected behavior, functional capabilities,
non-functional requirements, user interactions, data requirements, system
constraints, and testing requirements. The SRS will serve as the primary
reference for the design, development, testing, and evaluation of the
Dikaios Amanah system.

The system is designed around the Musharakah partnership model, where
investment participants share profits and losses according to mutually
agreed terms. The proposed platform emphasizes transparency, verification,
security, and traceability throughout the investment lifecycle.


\subsection{Scope}

Dikaios Amanah will provide a web-based environment through which verified
business owners can present eligible SME investment opportunities and
investors can review and participate in those opportunities.

The system will support three primary user roles:

\begin{itemize}
    \item Business Owner
    \item Investor
    \item Administrator
\end{itemize}

Business owners will be able to create and manage business profiles,
submit required verification information, upload supporting documents,
create investment campaigns after approval, provide business progress
information, and manage campaign-related information.

Investors will be able to create accounts, complete the required
verification process, browse approved investment campaigns, review
business and campaign information, participate in simulated investment
transactions, and monitor their investment portfolio and history.

Administrators will be responsible for platform-level management activities
including user verification, business approval, document review, campaign
moderation, risk assessment, user management, and monitoring of system
activity.

The system will also provide an audit and trust mechanism for important
system events. The proposed architecture includes a Permissioned Trust and
Audit Layer intended to improve transparency and provide tamper-evident
records of critical activities.

The Minimum Viable Product (MVP) will focus on the core platform workflow
required for academic demonstration and evaluation. Real banking
integration, live payment processing, production-level regulatory
implementation, AI-based fraud detection, AI-powered risk prediction,
production Hyperledger Fabric deployment, and automated smart-contract
execution are outside the scope of the MVP.


\subsection{Intended Audience}

This document is intended for the following audiences:

\begin{itemize}
    \item Project developer and system designer
    \item Software Engineering course instructor and evaluator
    \item Software testing and quality assurance personnel
    \item Project stakeholders
    \item Future developers who may extend the system
\end{itemize}

The document provides the requirements and technical specifications needed
to understand the expected behavior and structure of the proposed system.


\subsection{Product Overview}

Dikaios Amanah is intended to function as a technology-enabled platform
connecting verified SMEs seeking investment with verified investors seeking
structured partnership-based investment opportunities.

The proposed system follows the Musharakah principle rather than
interest-based lending. The platform is designed to provide a transparent
investment lifecycle beginning with registration and verification and
continuing through campaign approval, investment participation, agreement
management, business reporting, profit declaration, and profit
distribution.

The system architecture proposed in the project consists of a web
frontend, backend services, a relational database, secure document
storage, authentication and security mechanisms, and a Permissioned Trust
and Audit Layer.

The application layer is planned around React for frontend development and
Node.js with Express.js for backend services. PostgreSQL will be used for
relational data management, while Firebase Storage is intended for secure
document storage. Authentication and authorization mechanisms will protect
user accounts and role-specific functionality.

Critical events and records will be subject to audit mechanisms so that
important activities can be reviewed by authorized users. The trust and
audit mechanism is intended to complement the primary application database
rather than replace it.


\subsection{Definitions and Acronyms}

\begin{description}

    \item[SME] Small and Medium Enterprise.

    \item[MVP] Minimum Viable Product.

    \item[KYC] Know Your Customer.

    \item[RBAC] Role-Based Access Control.

    \item[JWT] JSON Web Token.

    \item[API] Application Programming Interface.

    \item[HTTPS] Hypertext Transfer Protocol Secure.

    \item[SHA-256] Secure Hash Algorithm 256-bit.

    \item[SRS] Software Requirements Specification.

    \item[UML] Unified Modeling Language.

    \item[SDLC] Software Development Life Cycle.

    \item[ERD] Entity Relationship Diagram.

\end{description}


\subsection{References}

The primary project-specific source for this Software Requirements
Specification is the Dikaios Amanah Project Proposal. The proposal
establishes the project's background, objectives, scope, proposed solution,
system architecture, technology stack, feasibility considerations,
limitations, and planned development activities.

Additional technical and domain references are listed in the References
section of this document.
